% Chapter 4, Topic _Linear Algebra_ Jim Hefferon
%  http://joshua.smcvt.edu/linearalgebra
%  2001-Jun-12
\topic{线性递推}
\index{linear recurrence|(}
\index{recurrence relation|(}

1202 年，以 Fibonacci（斐波那契）之名为人熟知的
比萨的 Leonardo 提出了这样一个问题。
\begin{quotation}
  \noindent 有人把一对兔子放进一个四周围墙的地方。
  假设每对兔子每个月生出一对新兔子，而新生的兔子
  从第二个月起便具有繁殖能力，
  那么一年之内由这一对兔子可以繁殖出多少对兔子？
\end{quotation}
它超出了初等的种群指数增长模型，
把新生个体有一段时间不具繁殖能力（这里是一个月）这一点也考虑进来。
不过，它保留了其他一些简化假设，比如
兔子在过了某个年龄之后便不再有繁殖能力。

要知道下个月的兔子对数总和，
我们把进入下个月时还活着的兔子对数，
加上下个月将新出生的兔子对数。
后者等于进入下个月时具有繁殖能力的兔子对数，
也就是到下个月时已经存活至少两个月的兔子对数。
\begin{equation*}
  F(n)=F(n-1)+F(n-2)  \qquad \text{其中 $F(0)=0$，$F(1)=1$}
  \tag{$*$}
\end{equation*}
左边是一个
\definend{递推关系}。\index{recurrence}
这个名称的由来是
$F$ 出现在它自己的定义式中。
右边是初始条件。
由~($*$) 我们可以计算 $F(2)$、$F(3)$ 等，
从而逐步得到 Fibonacci 问题的答案。
\begin{center}
  \begin{tabular}{r|ccccccccccccc}
    \textit{月份}~$n$
     &0  &1  &2  &3  &4  &5  &6  &7  &8  &9  &10  &11  &12  \\ \hline
    \textit{对数}~$F(n)$
     &0  &1  &1  &2  &3  &5  &8  &13  &21  &34  &55  &89  &144   
  \end{tabular}
\end{center}
我们将利用线性代数得到一个公式，
它无需先算出中间值 $F(2)$、$F(3)$ 等
就能计算 $F(n)$。

我们先把 ($*$) 写成矩阵形式。
\begin{equation*}
  \colvec{F(n) \\ F(n-1)}
  =
  \begin{mat}[r]
    1  &1   \\
    1  &0
  \end{mat}
  \colvec{F(n-1) \\ F(n-2)}
  \qquad\text{其中 $\colvec{F(1) \\ F(0)}=\colvec{1 \\  0}$}
\end{equation*}  
记这个矩阵为 $T$，
记分量为 $F(n)$ 与 $F(n-1)$ 的向量为 $\vec{v}_{n}$，
则当 $n\geq 1$ 时 $\vec{v}_{n}=T^{n-1}\vec{v}_1$。
如果将 $T$ 对角化，那么
我们就有了计算它的幂的快速方法：~若 $T=PDP^{-1}$，则
$T^n=PD^nP^{-1}$，而对角矩阵 $D$ 的 $n$ 次幂就是这样的对角矩阵，
其对角元是 $D$ 的对角元各自的 $n$ 次幂。

$T$ 的特征方程是 $\lambda^2-\lambda-1=0$。
求根公式给出它的两个根为 $(1+\sqrt{5})/2$ 与
$(1-\sqrt{5})/2$。
（它们有时被称为"黄金比"；
参见 \cite{Falbo}。）
对角化给出下式。
\begin{equation*}
  \begin{mat}[r]
    1  &1  \\
    1  &0
  \end{mat}
  =\begin{mat}
     \frac{1+\sqrt{5}}{2}  &\frac{1-\sqrt{5}}{2} \\
     % \sfrac{1+\sqrt{5}}{2}  &\sfrac{1-\sqrt{5}}{2} \\
     1                     &1
   \end{mat}
   \begin{mat}
     \frac{1+\sqrt{5}}{2}  &0   \\
     0                     &\frac{1-\sqrt{5}}{2}
   \end{mat}
   \begin{mat}
     \frac{1}{\sqrt{5}}     &-(\frac{1-\sqrt{5}}{2\sqrt{5}})  \\
     \frac{-1}{\sqrt{5}}    &\frac{1+\sqrt{5}}{2\sqrt{5}}       
   \end{mat}
\end{equation*} 
引入向量并取 $n$ 次幂，我们得到
\begin{multline*}
  \colvec{F(n) \\ F(n-1)}
  =\begin{mat}[r]
    1  &1  \\
    1  &0
  \end{mat}^{n-1}
  \colvec{f(1) \\ f(0)}                                          \\
  =\begin{mat}
     \frac{1+\sqrt{5}}{2}  &\frac{1-\sqrt{5}}{2} \\
     1                     &1
   \end{mat}
   \begin{mat}
     \left(\frac{1+\sqrt{5}}{2}\right)^{n-1}  &0   \\
     0                                  &\left(\frac{1-\sqrt{5}}{2}\right)^{n-1}
   \end{mat}
   \begin{mat}
     \frac{1}{\sqrt{5}}     &-(\frac{1-\sqrt{5}}{2\sqrt{5}})  \\
     \frac{-1}{\sqrt{5}}    &\frac{1+\sqrt{5}}{2\sqrt{5}}       
   \end{mat}
  \colvec[r]{1 \\ 0}                                             
\end{multline*}
这个计算虽然难看，但并不难。
\begin{align*}
  \colvec{F(n) \\ F(n-1)}
  &=\begin{mat}
     \frac{1+\sqrt{5}}{2}  &\frac{1-\sqrt{5}}{2} \\
     1                     &1
   \end{mat}
   \begin{mat}
     \left(\frac{1+\sqrt{5}}{2}\right)^{n-1}  &0   \\
     0                                  &\left(\frac{1-\sqrt{5}}{2}\right)^{n-1}
   \end{mat}
  \colvec{\frac{1}{\sqrt{5}} \\ -\frac{1}{\sqrt{5}}}       \\ 
  &=\frac{1}{\sqrt{5}}
    \begin{mat}
     \frac{1+\sqrt{5}}{2}  &\frac{1-\sqrt{5}}{2} \\
     1                     &1
   \end{mat}
  \colvec{\left(\frac{1+\sqrt{5}}{2}\right)^{n-1} \\ 
          -\left(\frac{1-\sqrt{5}}{2}\right)^{n-1}}    \\    
  &=\frac{1}{\sqrt{5}}
  \colvec{\left(\frac{1+\sqrt{5}}{2}\right)^{n}  
          -\left(\frac{1-\sqrt{5}}{2}\right)^{n}  \\
          \left(\frac{1+\sqrt{5}}{2}\right)^{n-1}  
          -\left(\frac{1-\sqrt{5}}{2}\right)^{n-1}}   
\end{align*}
我们要的是第一个分量。
\begin{equation*}
  F(n)=\frac{1}{\sqrt{5}}\left[\left(\frac{1+\sqrt{5}}{2}\right)^{n}
                               -\left(\frac{1-\sqrt{5}}{2}\right)^{n}\right]
\end{equation*}
这个公式给出序列中任何一项的值，
而不必先求出中间值。

由于 $(1-\sqrt{5})/2\approx -0.618$
的绝对值小于一，它的幂趋于零，所以
$F(n)$ 的公式由它的第一项主导。
尽管我们通过在开始具有繁殖能力之前
加入一段延迟期而扩展了初等的
种群增长模型，但我们
得到的仍是一个渐近指数式的函数。

一般地，一个
\definend{$k$ 阶齐次线性递推关系}\index{recurrence}\index{linear recurrence!definition} 
具有如下形式。
\begin{equation*}
  f(n)=a_{n-1}f(n-1)+a_{n-2}f(n-2)+\dots+a_{n-k}f(n-k)
\end{equation*}
这个递推关系是齐次的，
因为它没有常数项，也就是说，我们可以把它改写为
$0=-f(n)+a_{n-1}f(n-1)+a_{n-2}f(n-2)+\dots+a_{n-k}f(n-k)$。
它是~$k$ 阶的，
因为它用 $k$ 个此前的项来计算 $f(n)$。
该关系与给出
$f(0)$、\ldots、$f(k-1)$ 之值的
初始条件\index{recurrence!initial conditions} 
一起，完全确定一个序列，
原因很简单：
我们可以先计算 $f(k)$、$f(k+1)$ 等，再算出 $f(n)$。
与 Fibonacci 的情形一样，我们将求得一个求解该递推的公式，
它直接给出 $f(n)$。

设 $V$ 是定义域为
$\N=\set{0,1,2,\ldots}$、陪域为 $\C$ 的函数的集合。
（在方便时，我们有时使用定义域 $\Z^+=\set{1,2,\ldots}$。）
在通常的加法与
标量乘法的意义下，即
$f+g$ 是映射 $x\mapsto f(x)+g(x)$，
$cf$ 是映射 $x\mapsto c\cdot f(x)$，
这是一个向量空间。

如果我们撇开一切初始条件，只考察递推本身，
那么可能有许多函数满足该关系。
例如，Fibonacci 递推——除初始值之外的每个值
都是前两项之和——就为函数 $L$ 所满足，其前几个值
为 $L(0)=2$、$L(1)=1$、$L(2)=3$、$L(3)=4$，
以及 $L(4)=7$。

取定一个~$k$ 阶齐次线性递推关系，
考虑满足该关系（不带
初始条件）的函数组成的子集 $S$。
这个 $S$ 是 $V$ 的一个子空间。
它非空，因为由齐次性，零函数是一个解。
它对加法封闭，因为若 $f_1$ 与 $f_2$ 是解，则
该式成立。
\begin{multline*}
  -(f_1+f_2)(n)+a_{n-1}(f_1+f_2)(n-1)+\dots+a_{n-k}(f_1+f_2)(n-k) \\  
  \begin{aligned}
    &=(-f_1(n)+\dots+a_{n-k}f_1(n-k))          \\
    &\qquad\hbox{}+(-f_2(n)+\dots+a_{n-k}f_2(n-k))     \\
    &=0+0=0
  \end{aligned}
\end{multline*}
它对标量乘法也封闭。
\begin{multline*}
  -(rf_1)(n)+a_{n-1}(rf_1)(n-1)+\dots+a_{n-k}(rf_1)(n-k) \\  
  \begin{aligned}
    &=r\cdot (-f_1(n)+\dots+a_{n-k}f_1(n-k))   \\
    &=r\cdot 0                                    \\
    &=0
  \end{aligned}
\end{multline*}
我们可以求出 $S$ 的维数。
设 $k$ 是递推的阶，
考虑这个从函数集合~$S$ 到
$k$ 维高的向量集合的映射。
\begin{equation*}
  f 
  \mapsto 
  \colvec{f(0) \\ f(1) \\ \vdots \\ f(k-1)}
\end{equation*}
\nearbyexercise{exer:SeqToRnLinMap} 证明了这个映射是线性的。
递推的任何一个解都由那 $k$ 个
初始条件唯一确定，所以这个映射是单射且满射。
因此它是一个同构，且 $S$ 的维数为 $k$。

于是我们可以这样来描述~$k$ 阶线性齐次递推关系的
解集：求出一个由
$k$ 个线性无关的函数组成的基。
为了造出这些函数，我们把
递推写成矩阵形式。
\begin{equation*}
  \colvec{f(n) \\ f(n-1) \\ \vdots  \\ f(n-k+1)}
  =
  \begin{mat}
    a_{n-1}  &a_{n-2}  &a_{n-3}  &\ldots  &a_{n-k+1} &a_{n-k}  \\
    1    &0        &0        &\ldots  &0         &0        \\
    0    &1        &0                                      \\
    0    &0        &1                                      \\
    \vdots &\vdots &         &\ddots  &           &\vdots  \\
    0    &0        &0        &\ldots   &1          &0
  \end{mat}
  \colvec{f(n-1) \\ f(n-2) \\ \vdots  \\ f(n-k)}
\end{equation*}
将这个矩阵记为~$A$。
我们要求它的特征函数，
即 $A-\lambda I$ 的行列式。
$\nbyn{2}$ 情形中的模式
\begin{equation*}
  \begin{mat}
    a_{n-1}-\lambda  &a_{n-2} \\
    1               &-\lambda
  \end{mat}
  =\lambda^2-a_{n-1}\lambda-a_{n-2}
\end{equation*}
以及 $\nbyn{3}$ 情形
\begin{equation*}
  \begin{mat}
    a_{n-1}-\lambda  &a_{n-2}   &a_{n-3}  \\
    1               &-\lambda  &0        \\
    0               &1         &-\lambda
  \end{mat}
  =-\lambda^3+a_{n-1}\lambda^2+a_{n-2}\lambda+a_{n-3}
\end{equation*}
使我们预期，并且 \nearbyexercise{exer:CharEqnIsDeter} 也验证了：
这就是特征方程。
\begin{align*}
  0 &=
  \begin{vmat}
    a_{n-1}-\lambda &a_{n-2}  &a_{n-3}  &\ldots  &a_{n-k+1} &a_{n-k}  \\
    1              &-\lambda &0        &\ldots  &0         &0        \\
    0              &1        &-\lambda                                      \\
    0              &0        &1                                      \\
    \vdots         &\vdots &         &\ddots   &           &\vdots  \\
    0             &0        &0        &\ldots   &1          &-\lambda
  \end{vmat}                                                           \\
  &=\pm(-\lambda^k+a_{n-1}\lambda^{k-1}+a_{n-2}\lambda^{k-2}
       +\dots+a_{n-k+1}\lambda+a_{n-k})
\end{align*}
这个 $\pm$ 对求根没有影响，
所以我们把它去掉。
我们称多项式
$-\lambda^k+a_{n-1}\lambda^{k-1}+a_{n-2}\lambda^{k-2}
       +\dots+a_{n-k+1}\lambda+a_{n-k}$
与该递推关系相伴。\index{recurrence!associated polynomial}\index{polynomial!associated with recurrence}

如果特征方程没有重根，
那么该矩阵可对角化，
于是原则上我们可以像 Fibonacci 情形那样得到 $f(n)$ 的公式。
但由于我们知道解子空间的维数为~$k$，
因此无需做对角化计算，只要
我们能展示出 $k$ 个满足该关系的不同的
线性无关函数。

设 $r_1$、$r_2$、\ldots、$r_k$ 是互不相同的根，
考虑以这些根的幂为值的函数，即从 $f_{r_1}(n)=r_1^n$ 到
$f_{r_k}(n)=r_k^n$。
\nearbyexercise{exer:SoltnsLinRecur} 证明
每一个都是递推的解，并且它们组成一个
线性无关集。
% So given the homogeneous linear recurrence
% $f(n+1)=a_nf(n)+\dots+a_{n-k}f(n-k)$
% (that is, $0=-f(n+1)+a_nf(n)+\dots+a_{n-k}f(n-k)$)
% we consider the associated equation
% $0=-\lambda^k+a_n\lambda^{k-1}+\dots+a_{n-k+1}\lambda+a_{n-k}$.
于是，如果相伴多项式的
根互不相同，
那么该关系的任何解都具有形式
$f(n)=c_1r_1^n+c_2r_2^n+\dots+c_kr_k^n$，其中 $c_1, \dots, c_n$ 是一些标量。
（重根的情形类似，但我们
在这里不讨论；参见任何一本离散数学教材。）

现在把初始条件引入。
用它们解出 $c_1$、\ldots、$c_n$。
例如，与 Fibonacci 关系相伴的多项式是
$-\lambda^2+\lambda+1$，它的根是 $r_1=(1+\sqrt{5})/2$
和 $r_2=(1-\sqrt{5})/2$，
于是 Fibonacci 递推的任何解
都具有形式 $f(n)=c_1((1+\sqrt{5})/2)^n+c_2((1-\sqrt{5})/2)^n$。
取 $n=0$ 与 $n=1$ 时的 Fibonacci 初始条件
\begin{equation*}
  \begin{linsys}{2}
    c_1               &+  &c_2               &=  &0  \\
    (1+\sqrt{5}/2)c_1 &+  &(1-\sqrt{5}/2)c_2 &=  &1
  \end{linsys}
\end{equation*}
解得 $c_1=1/\sqrt{5}$ 与 $c_2=-1/\sqrt{5}$，正如我们在上面求得的。

最后我们考虑非齐次情形，
其中关系具有形式
$f(n+1)=a_nf(n)+a_{n-1}f(n-1)+\dots+a_{n-k}f(n-k)+b$，$b$ 为某个非零数。
从齐次情形过渡过来，
我们只需作一点小小的调整。

这个经典例子做了说明：
1883 年，Edouard Lucas 提出了
Hanoi 塔\index{Tower of Hanoi}问题。
\begin{quotation}
  \noindent
  在贝拿勒斯的大圣庙里，在标志世界中心的穹顶之下，
  安放着一块铜板，铜板上固定着三根金刚石针，
  每根高一腕尺，粗如蜜蜂的躯体。
  开天辟地之时，
  上帝在其中一根针上安放了六十四片纯金圆盘，最大的圆盘
  放在铜板上，其余的圆盘越来越小，
  直至最上面的一片。
  这就是梵塔。
  祭司们日夜不停地按照梵天一成不变的法则
  把圆盘从一根金刚石针移到另一根上，
  该法则要求值班的祭司一次
  不得移动多于一片圆盘，
  并且必须把这片圆盘放在一根针上，使得它下面
  没有更小的圆盘。
  当这六十四片圆盘就这样从创世时
  上帝安放它们的那根针被移到另一根针上时，
  塔、庙宇和婆罗门都将化为尘土，
  随着一声霹雳，世界也将消失。
  （\cite{DeParville} 的译文，转引自 \cite{Ball}。）
\end{quotation}
我们撇开"祭司们为什么不坐下歇一会儿、
好让世界多存在一段时间"这个问题不谈，转而问
需要移动圆盘多少次。
在着手解决六十四片圆盘的问题之前，
我们先考虑
三片圆盘的问题。

一开始，三片圆盘都在同一根针上。
\begin{center}
  \includegraphics{jc/mp/ch5.6}
\end{center}
经过三次移动——把小圆盘移到远端的针上，
把中等大小的圆盘移到
中间的针上，再把小圆盘移到中间的针上——
我们就得到下图。
\begin{center}
  \includegraphics{jc/mp/ch5.7}
\end{center}
现在我们可以把大圆盘移到远端的针上。
然后，为完成整个移动，我们对两片较小的
圆盘重复这三步移动过程，这一次要让它们最终
落在第三根针上、
大圆盘的上面。

这一移动序列是我们所能做到的最好结果。
要以最少的次数移动最底下的圆盘，
我们必须先把较小的圆盘移到中间的针上，
然后移动大圆盘，
最后把所有较小的圆盘从中间的针
移到目标针上。
由于这个最少次数就够了，我们得到如下递推。
\begin{equation*}
  T(n)=T(n-1)+1+T(n-1)=2T(n-1)+1 \qquad \text{其中 $T(1)=1$}
\end{equation*} 
下面是 $T$ 的前几个值。
\begin{center}
  \begin{tabular}{r|cccccccccc}
    \textit{圆盘数}~$n$             
       &1  &2  &3  &4  &5  &6     &7    &8    &9   &10  \\
    \hline
    \textit{移动次数}~$T(n)$ 
       &1  &3  &7  &15  &31  &63  &127  &255  &511 &1023 
  \end{tabular}
\end{center}
当然，这些数都比 2 的幂小一。
为推导这一点，
把原关系写成 $-1=-T(n)+2T(n-1)$。
考虑 $0=-T(n)+2T(n-1)$，一个~$1$ 阶的线性齐次递推。
它的相伴多项式是 $-\lambda+2$，
有唯一的根 $r_1=2$。
于是
满足该齐次关系的函数具有形式 $c_12^n$。

这就是齐次解。
现在我们需要一个特解。
由于非齐次关系 $-1=-T(n)+2T(n-1)$ 非常简单，
我们可以一眼看出一个特解 $T(n)=-1$。
递推
$T(n)=2T(n-1)+1$（不带初始条件）
的任何解都是齐次解与
特解之和：$c_12^n-1$。
现在，初始条件 $T(1)=1$ 给出 $c_1=1$，于是我们得到了
生成上表的公式：~$n$ 片圆盘的 Hanoi 塔问题
需要 $T(n)=2^n-1$ 次移动。

在更复杂的情形下求特解，
也许并不奇怪，
也更加复杂。
一份令人愉快且颇有收获、但具有挑战性的资料
是 \cite{GrahamKnuthPatashnik}。
关于 Hanoi 塔的更多内容，参见 \cite{Ball}、\cite{Gardner57}，
以及 \cite{Hofstadter}。
练习之后附有一些
计算机代码。

\begin{exercises}
  \item 
   多少个月后 Fibonacci 兔子对的数量
   会超过一千？
   一万呢？
   一百万呢？
   \begin{answer}
     我们使用公式
     \begin{equation*}
       F(n)=\frac{1}{\sqrt{5}}\left[\left(\frac{1+\sqrt{5}}{2}\right)^{n}
                               -\left(\frac{1-\sqrt{5}}{2}\right)^{n}\right]
     \end{equation*}
     如前面所观察到的，$(1+\sqrt{5})/2$ 大于一，
     而 $(1+\sqrt{5})/2$ 的绝对值小于一。
\begin{lstlisting}
sage: phi = (1+5^(0.5))/2
sage: psi = (1-5^(0.5))/2
sage: phi
1.61803398874989
sage: psi
-0.618033988749895       
\end{lstlisting}
     所以该表达式的值由第一项主导。
     解 $1000=(1/\sqrt{5})\cdot ((1+\sqrt{5})/2)^n$ 给出
     下式。
\begin{lstlisting}
sage: a = ln(1000*5^(0.5))/ln(phi)
sage: a
16.0271918385296
sage: psi^(17)
-0.000280033582072583       
\end{lstlisting}
   于是到十七次幂时，第二项的贡献已不足以
   改变舍入的结果。
   对于一万和一百万的计算，情况
   更加极端。
\begin{lstlisting}
sage: b = ln(10000*5^(0.5))/ln(phi)
sage: b
20.8121638053112
sage: c = ln(1000000*5^(0.5))/ln(phi)
sage: c
30.3821077388746    
\end{lstlisting}
   在这两种情形下，答案是 $21$ 和 $31$。
   \end{answer}
  \item 
   求解下列各个齐次线性递推关系。
   \begin{exparts}
     \partsitem $f(n)=5f(n-1)-6f(n-2)$   
     \partsitem $f(n)=4f(n-2)$   
     \partsitem $f(n)=5f(n-1)-2f(n-2)-8f(n-3)$   
   \end{exparts}
   \begin{answer}
    \begin{exparts}
      \partsitem 
        我们把该关系写成矩阵形式。
        \begin{equation*}
          \colvec{f(n) \\ f(n-1)}
          =
          \begin{mat}[r]
            5  &-6  \\
            1  &0
          \end{mat}
          \colvec{f(n-1) \\ f(n-2)}
        \end{equation*}
        这个矩阵的特征方程
        \begin{equation*}
          \begin{vmat}
            5-\lambda &-6       \\
            1         &-\lambda
          \end{vmat}
          =\lambda^2-5\lambda+6 
        \end{equation*}
        有根 $2$ 和 $3$。
        任何形如
        $f(n)=c_12^n+c_23^n$
        的函数都满足该递推。
      \partsitem 
        该关系的矩阵表示是
        \begin{equation*}
          \colvec{f(n) \\ f(n-1)}
          =
          \begin{mat}[r]
            0  &4  \\
            1  &0  
          \end{mat}
          \colvec{f(n-1) \\ f(n-2)}
        \end{equation*}
        而特征方程
        \begin{equation*}
          \begin{vmat}
            \lambda^2-2       
          \end{vmat}
          =(\lambda-2)(\lambda+2)
        \end{equation*}
        有两个根 $2$ 和 $-2$。
        任何形如
        $f(n)=c_12^n+c_2(-2)^n$
        的函数都满足这个递推。
      \partsitem 
        用矩阵形式，该关系
        \begin{equation*}
          \colvec{f(n) \\ f(n-1) \\ f(n-2)}
          =
          \begin{mat}[r]
            5  &-2  &-8  \\
            1  &0  &0  \\
            0  &1  &0
          \end{mat}
          \colvec{f(n-1) \\ f(n-2) \\ f(n-3)}
        \end{equation*}
        的特征方程有根 $-1$、$2$ 和 $4$。
        任何形如
        $c_1(-1)^n+c_22^n+c_34^n$ 的组合都是该递推的解。
    \end{exparts} 
   \end{answer}
  \item \label{exer:SolvePartRecurSoltn} 
    对前一练习中的各关系，给出带
    下列初始条件的公式。
    \begin{exparts}
      \partsitem $f(0)=1$, $f(1)=1$
      \partsitem $f(0)=0$, $f(1)=1$
      \partsitem $f(0)=1$, $f(1)=1$, $f(2)=3$.      
    \end{exparts}
    \begin{answer}
      \begin{exparts}
        \partsitem 该齐次递推的解是
          $f(n)=c_12^n+c_23^n$。
          代入 $f(0)=1$ 与 $f(1)=1$ 给出这个线性方程组。
          \begin{equation*}
            \begin{linsys}{2}
              c_1 &+ &c_2   &= &1 \\
             2c_1 &+ &3c_2  &= &1
            \end{linsys}
          \end{equation*}
          一眼可见 $c_1=2$，$c_2=-1$。
        \partsitem
          该齐次递推的解是
          $c_12^n+c_2(-2)^n$。
          初始条件给出这个线性方程组。
          \begin{equation*}
            \begin{linsys}{2}
              c_1 &+ &c_2  &= &0  \\
             2c_1 &- &2c_2 &= &1
            \end{linsys}
          \end{equation*}
          解为 $c_1=1/4$，$c_2=-1/4$。
        \partsitem
          该齐次递推有解
          $f(n)=c_1(-1)^n+c_22^n+c_34^n$。
          结合初始条件，我们得到这个线性方程组。
          \begin{equation*}
            \begin{linsys}{3}
              c_1 &+  &c_2   &+  &c_3  &= &1 \\
             -c_1 &+  &2c_2  &+  &4c_3  &= &1 \\
              c_1 &+  &4c_2  &+  &16c_3  &= &3
            \end{linsys}
          \end{equation*}
          其解为 $c_1=1/3$，$c_2=2/3$，$c_3=0$。
      \end{exparts}
    \end{answer}
  \item \label{exer:SeqToRnLinMap}
    验证在 $S$ 与 $\Re^k$ 之间给出的那个同构是线性映射。
    % We argue above that this map is one-to-one.
    % What is its inverse?
    \begin{answer}
      取定一个~$k$ 阶线性齐次递推。
      设 $S$ 是满足该递推的
      函数 $\map{f}{\N}{\C}$ 的集合。
      考虑如下给出的函数 $\map{\Phi}{S}{\C^k}$。
      \begin{equation*}
        f\mapsunder{\Phi}\colvec{f(0) \\ f(1) \\ \vdots \\ f(k-1)}
      \end{equation*}
      这就证明了线性性。
      \begin{multline*}
        \Phi(a\cdot f_1+b\cdot f_2)
        =
        \colvec{af_1(0)+bf_2(0) \\ \vdots \\ af_1(k-1)+bf_2(k-1)}     \\
        =
        a\colvec{f_1(0) \\ \vdots \\ f_1(k-1)}
        +
        b\colvec{f_2(0) \\ \vdots \\ f_2(k-1)}
        =a\cdot \Phi(f_1)+b\cdot \Phi(f_2)
      \end{multline*}
    \end{answer}
  \item \label{exer:CharEqnIsDeter}
    证明该矩阵的特征方程正如所述，也就是说，
    它是与该关系相伴的多项式。
    （\textit{提示：}沿最后一列展开
    并使用归纳法即可。）
    \begin{answer}
      我们利用提示来证明这一点。
      \begin{align*}
        0 &=
        \begin{vmat}
          a_{n-1}-\lambda &a_{n-2}  &a_{n-3}  &\ldots  &a_{n-k+1} &a_{n-k}  \\
          1              &-\lambda &0        &\ldots  &0         &0        \\
          0              &1        &-\lambda                               \\
          0              &0        &1                                      \\
          \vdots         &\vdots &         &\ddots   &           &\vdots  \\
          0             &0        &0        &\ldots   &1          &-\lambda
        \end{vmat}                                                           \\
        &=(-1)^{k-1}(-\lambda^k+a_{n-1}\lambda^{k-1}+a_{n-2}\lambda^{k-2}       
               +\dots+a_{n-k+1}\lambda+a_{n-k})                      \\
        &=\pm(-\lambda^k+a_{n-1}\lambda^{k-1}+a_{n-2}\lambda^{k-2}       
               +\dots+a_{n-k+1}\lambda+a_{n-k})
      \end{align*}
      基础步骤是平凡的。
      对于归纳步骤，
      沿最后一列展开给出两个非零项。
      \begin{equation*}
        (-1)^{k-1}a_{n-k}\cdot 1
        -\lambda\cdot
        \begin{vmat}
          a_{n-1}-\lambda &a_{n-2}  &a_{n-3}  &\ldots  &a_{n-k+1}  \\
          1              &-\lambda &0        &\ldots  &0       \\
          0              &1        &-\lambda                               \\
          0              &0        &1                                      \\
          \vdots         &\vdots &         &\ddots   &        \\
          0             &0        &0        &\ldots   &1     
        \end{vmat}                                                          
      \end{equation*}
      （该矩阵是方阵，所以 $-\lambda$
      前面的符号是 $-1^{\text{偶}}$）。
      应用归纳假设就得到想要的结果。
      \begin{multline*}
        =(-1)^{k-1}a_{n-k}\cdot 1                     \\
        -\lambda\cdot
         (-1)^{k-2}(-\lambda^{k-1}+a_{n-1}\lambda^{k-2}+a_{n-2}\lambda^{k-3}       
               +\dots+a_{n-k+1}\lambda^0)
      \end{multline*}
    \end{answer}
  \item \label{exer:SoltnsLinRecur}
    给定齐次线性递推关系
    $f(n)=a_nf(n-1)+\dots+a_{n-k}f(n-k)$，设 $r_1$、\ldots、$r_k$ 是
    相伴多项式的根。
    证明每个函数
         $f_{r_i}(n)=r_k^n$
         都满足该递推（不带初始条件）。
    \begin{answer}
       这是对 $n$ 的直接归纳。
    \end{answer}
  \item 
    （这指的是下面计算机代码中给出的值
    $T(64)=18,446,744,073,709,551,615$。）
    若每秒移动一片圆盘，
    Hanoi 塔的祭司们需要多少年才能完成这项工作？
    \begin{answer}
      \Sage{} 表明我们是安全的。
\begin{lstlisting}
sage: T64 = 18446744073709551615
sage: T64_days = T64/(60*60*24)
sage: T64_days
1229782938247303441/5760
sage: T64_years = T64_days/365.25
sage: T64_years
5.84542046090626e11        
sage: age_of_universe = 13.8e9
sage: T64_years/age_of_universe
42.3581192819294
\end{lstlisting}
    \end{answer}

\end{exercises}

\announcecomputercode
这段代码生成由一个递推关系和初始条件
所定义的函数的前几个值。
它用的是 LISP 的 Scheme 方言，
具体是 \cite{ChickenScheme}。

在加载了一个扩展——它使计算机在整数变大时
不会切换到浮点数——之后，
Hanoi 塔函数就很简单直白了。
\begin{lstlisting}
(require-extension numbers)

(define (tower-of-hanoi-moves n) 
    (if (= n 1)
       1
       (+ (* (tower-of-hanoi-moves (- n 1)) 
              2) 
           1) ) )

; Two helper funcitons
(define (first-few-outputs proc n)
    (first-few-outputs-aux proc n '()) )

(define (first-few-outputs-aux proc n lst)
    (if (< n 1)
    lst 
    (first-few-outputs-aux proc (- n 1) (cons (proc n) lst)) ) )
\end{lstlisting}
\noindent （对于不熟悉递归代码的读者：~为计算 $T(64)$，
计算机想要计算 $2*T(63)-1$，这需要
先算出 $T(63)$。
计算机把"乘 $2$"和"加 $1$"暂时放在一边。
它用这同一段代码来计算 $T(63)$（这就是
"递归"的含义），而为此它想计算 $2*T(62)-1$。
如此继续，经过 $63$ 步之后，计算机
尝试计算 $T(1)$。
然后它返回 $T(1)=1$，
这使得 $T(2)$ 的计算得以进行，
如此等等，直到最初对 $T(64)$ 的计算完成。）

这些辅助函数给出前几个
值的表。
% (Some language notes:~\texttt{'()} is the empty list
% and \texttt{cons} pushes something onto the start of a 
% list.
% Note that, in the last line, the function \texttt{proc}
% is called on argument \texttt{n}.)
下面是提示符处的会话。
\begin{lstlisting}
#;1> (load "hanoi.scm")
; loading hanoi.scm ...
; loading /var/lib//chicken/6/numbers.import.so ...
; loading /var/lib//chicken/6/chicken.import.so ...
; loading /var/lib//chicken/6/foreign.import.so ...
; loading /var/lib//chicken/6/numbers.so ...
#;2> (tower-of-hanoi-moves 64)
18446744073709551615
#;3> (first-few-outputs tower-of-hanoi-moves 64)
(1 3 7 15 31 63 127 255 511 1023 2047 4095 8191 16383 32767 65535 131071 262143 524287 1048575 
2097151 4194303 8388607 16777215 33554431 67108863 134217727 268435455 536870911 1073741823 
2147483647 4294967295 8589934591 17179869183 34359738367 68719476735 137438953471 274877906943 
549755813887 1099511627775 2199023255551 4398046511103 8796093022207 17592186044415 
35184372088831 70368744177663 140737488355327 281474976710655 562949953421311 1125899906842623 
2251799813685247 4503599627370495 9007199254740991 18014398509481983 36028797018963967 
72057594037927935 144115188075855871 288230376151711743 576460752303423487 1152921504606846975 
2305843009213693951 4611686018427387903 9223372036854775807 18446744073709551615)
\end{lstlisting}
\noindent 这是从 $T(1)$ 到 $T(64)$ 的列表（该会话经
编辑插入了换行以便阅读）。
\index{recurrence relation|)}
\index{linear recurrence|)}
\endinput
